\section{Integral descent and the Jacobian in a formal kernel sum}
\label{sec:geometry-formal-trace}

The denominator in a sum over a formal kernel comes from its differential
Jacobian.  It need not be the order of that kernel.  The following argument
keeps these two quantities separate and proves integral descent before
taking a trace.

Let $R$ be a complete discrete valuation ring of characteristic zero,
with maximal ideal $\mathfrak m$, residue field $k$, and fraction field
$K$.  Let $G$ be a smooth commutative formal group of dimension $r$ over
$R$, with chosen coordinates.  A small point is a tuple in $\mathfrak m^r$.
For such a point $a$, translation defines an $R$-algebra automorphism
\[
            t_a^*f(x)=f(x+_G a)\quad\hbox{of }R[[x]].
\]
The substitution converges in the maximal-ideal topology; its inverse is
translation by $-_G a$.

Let $p,q,s$ be endomorphisms of $G$ and put $\mu=p\circ q$.
Assume $s\circ\mu=[N]$ for an integer $N\geq1$.  Assume that the
coordinates of $[N]$ have finite colength in $k[[x]]$ after reduction.
Finally, assume that every geometric generic point of the kernel of
$\mu$ has coordinates in $R$.  This last condition is the splitting
hypothesis; finite height and splitting are retained geometric inputs.
Write $j_p=\det(dp_0)\in R$.

\begin{theorem}\label{thm:formal-kernel-jacobian-trace}
The maps $p,q,\mu$ are finite and flat on the corresponding power-series
rings.  Their differential determinants are nonzero.  Let
$d_q=\deg(q)$.  The small kernels of $q$ and $\mu$ are finite, and their
orders equal the respective degrees.  If $f\in R[[x]]$ satisfies
$t_a^*f=f$ for all $a\in\ker(q)$, there is a $g\in R[[x]]$ such that
\begin{equation}\label{eq:formal-kernel-exact-factor}
       \sum_{a\in\ker(\mu)}t_a^*f
                    =d_q\,j_p\,\mu^*g.
\end{equation}
In particular the factor in this identity is
$|\ker(q)|\det(dp_0)$, with its original normalization.
\end{theorem}

\begin{proof}
\emph{Finiteness and flatness.}
Since $s(0)=0$, each coordinate of $s\circ\mu$ belongs to the ideal
generated by the coordinates of $\mu$.  Therefore these coordinates have
finite colength after reduction.  Set
\[
 A_0=R[[z]],\qquad A_1=R[[y]],\qquad B=R[[x]],
       \qquad A_0\xrightarrow{p^*}A_1\xrightarrow{q^*}B.
\]
The composite is $\mu^*=q^*p^*$.  The ideal
$(\mathfrak m,\mu_1,\ldots,\mu_r)$ has radical the maximal ideal of $B$.
Choose lifts of a finite $k$-basis of its quotient.  Successive division
modulo powers of this ideal expresses every element of $B$ as an
$A_0$-linear combination of those lifts: the coefficients converge in
$A_0$, and the remainders tend to zero.  Hence $B$ is finite over $A_0$.

Differentiating $s\circ\mu=[N]$ gives
$ds_0\,dp_0\,dq_0=N I_r$.  All three determinants are nonzero, and the
formal inverse theorem over $K$ shows that $p^*,q^*,\mu^*$ are injective.
We may consequently regard the displayed rings as inclusions.
The same generators show that $B$ is finite over $A_1$; moreover $A_1$
is finite over $A_0$, being a submodule of the finite $A_0$-module $B$.
Each ring is regular local of dimension $r+1$, and each closed fibre has
dimension zero.  Miracle flatness therefore makes both maps flat
\cite{stacks-formal-trace-flatness}.  Their finite modules are free.

\emph{The different and its exact constant.}
For a formal group law $F(x,y)$ put
$C(x)=\partial F(x,y)/\partial y|_{y=0}$.  It is an integral matrix with
$C(0)=I_r$.  Differentiating the homomorphism identity for $p$ gives
\[
                 Dp(x)\,C(x)=C(p(x))\,dp_0.
\]
It follows that
\begin{equation}\label{eq:formal-jacobian-unit}
        \det Dp(x)=j_p\frac{\det C(p(x))}{\det C(x)}
                                      \in j_p R[[x]]^\times.
\end{equation}
The finite flat map $A_0\to A_1$ is a complete intersection.
Indeed its closed fibre is the Artinian ring
$k[[y]]/(\bar p_1,\ldots,\bar p_r)$.  The $r$ parameters form a regular
sequence in the regular ring $k[[y]]$.  The complete-intersection criterion
for the closed fibre of a flat finitely presented map gives the assertion
\cite{stacks-formal-trace-syntomic}.
This is a statement about the finite algebra $A_1$ over $A_0$; an
algebraization of the formal group is unnecessary.

For completeness, its ordinary relative differential module is the finite
module $\operatorname{coker}(Dp(y)^{t})$.  To justify using formal
derivatives here, every $A_0$-linear derivation to a finite $A_1$-module
is continuous in the $A_0$-adic topology.  That topology agrees with the
maximal-ideal topology because the closed fibre is Artinian.  The chain
rule therefore follows by taking limits of polynomials.  Conversely the
same convergent formula defines a derivation for every tuple killed by
the Jacobian relations.  This computes the ordinary differential module
of the \emph{finite algebra}, rather than identifying all algebraic
differentials of an arbitrary power-series ring with continuous ones.
The complete-intersection different formula now identifies the different
with its zeroth Fitting ideal, namely the Jacobian ideal
\cite{stacks-formal-trace-different}.  Equation~\eqref{eq:formal-jacobian-unit}
therefore gives
\[
 \mathfrak D_{A_1/A_0}=j_p A_1,\qquad
 \operatorname{Hom}_{A_0}(A_1,A_0)=j_p^{-1}A_1
\]
under the generic trace pairing.  The second equality says concretely
that $b\mapsto\operatorname{Tr}_{A_1/A_0}(b/j_p)$ is integral for every
$b\in A_1$.  Since $j_p$ is a constant in the base ring,
\begin{equation}\label{eq:formal-trace-jacobian-divisibility}
                   \operatorname{Tr}_{A_1/A_0}(A_1)\subseteq j_p A_0.
\end{equation}

\emph{Splitting and integral descent.}
The finite kernel algebra of $\mu$ is free over $R$, of rank $\deg\mu$.
Its generic fibre is reduced: its Jacobian is a nonzero constant times a
unit by the preceding calculation.  Hence it has $\deg\mu$ geometric
points.  By splitting their coordinates lie in $R$, and in fact in
$\mathfrak m$.  To see the latter assertion, if a coordinate had value
$c\in R^\times$, the unit $1-c^{-1}x_i$ of $R[[x]]$ would evaluate to
zero, which is impossible.  Thus these points give exactly $\deg\mu$
distinct integral translation automorphisms of $B$ over $A_0$.
Evaluation at the small coordinates agrees with the generic point by
continuity in the finite kernel algebra.  The same argument applies to
$q$, since its kernel is contained in that of $\mu$.

The fraction-field extension $\operatorname{Frac}(B)/\operatorname{Frac}(A_1)$
has degree $d_q$ and has $d_q$ distinct automorphisms given by these
translations.  It is therefore Galois with precisely this group.
The invariance of $f$ places it in $\operatorname{Frac}(A_1)$.
It is also integral over $A_1$, because $B$ is finite over $A_1$.
Normality of the regular ring $A_1$ proves $f\in A_1$; in the original
coordinates this means $f=q^*f_0$ for an \emph{integral} series $f_0$.
No division by a kernel order occurs in this descent.

\emph{The trace in the tower.}
The split translations identify the sum on the left of
\eqref{eq:formal-kernel-exact-factor} with the pullback of
$\operatorname{Tr}_{B/A_0}(f)$.
As $f=q^*f_0$, transitivity of trace gives
\[
 \operatorname{Tr}_{B/A_0}(f)
          =d_q\operatorname{Tr}_{A_1/A_0}(f_0)=d_qj_p g
                  \quad\hbox{for some }g\in A_0,
\]
using \eqref{eq:formal-trace-jacobian-divisibility}.  Pulling back through
$\mu$ proves the result.
\end{proof}

\paragraph{An integral gauge does not change the factor.}
Suppose $\kappa\in R[[x]]^\times$ has constant coefficient $1$, and the
given translation factors satisfy
\[
              U_a\,t_a^*\kappa=\kappa
                              \quad(a\in\ker\mu).
\]
If $U_a t_a^*f=f$ for $a\in\ker q$, then $\kappa^{-1}f$ is invariant
under $\ker q$.  The theorem, multiplied by $\kappa$, gives
\begin{equation}\label{eq:formal-kernel-twisted-factor}
       \sum_{a\in\ker\mu}U_a t_a^*f
                      =d_qj_p\,\kappa\mu^*g,
                         \qquad g\in R[[x]].
\end{equation}
Here $U_a$ and $\kappa$ are the original factors.  In particular a
rescaling of the section is unnecessary.

\paragraph{The canonical gauge in the theta application.}
Take the product formal group $G^n$ with its given integral
$O$-linear biextension, logarithmic gauge $h$, and $h_L$ as in
Theorem~\ref{thm:discriminant-harmonic-clearing}.  Let $m$ be a positive
odd integer, and let the matrix $\mu$
satisfy $\bar\mu^{t}\mu=m I_n$.
Then the precise factor is
\begin{equation}\label{eq:formal-kernel-canonical-kappa}
 \kappa=\exp\left(\sum_b h_L((\mu x)_b)
                              -m\sum_b h_L(x_b)\right)\in R[[x]]^\times,
                         \qquad\kappa(0)=1.
\end{equation}
Its integrality has a direct biextension proof.  Work modulo the additive
group of series with zero constant coefficient whose exponential is an
integral unit.  Substitution by integral formal points with zero constant
term preserves this group.
The biextension identities make $h$ additive in each argument modulo this
group, and the adjoint identity permits moving $a$ from the first argument
to $\bar a$ in the second.  Expanding the left side of
\eqref{eq:formal-kernel-canonical-kappa} before exponentiating gives
\[
 \sum_{i,j,k}h(x_j,\bar\mu_{ij}w\mu_{ik}x_k)
                      \equiv m\sum_jh(x_j,w x_j),
\]
which proves the assertion by the similitude equation.  All points in this
calculation have zero constant coefficient.

The supplied symmetric translation lifts obey
$U_a t_a^*\kappa=\kappa$.  Here is the normalization check.
Set $\chi=U_a(t_a^*\kappa)/\kappa$.
The cube identities give $t_a^*\chi=c\chi$ with $c$ a constant
unit.  The order identity is
$\prod_{i=0}^{m-1}(t_a^*)^i\chi=1$, hence
$c^{m(m-1)/2}\chi^m=1$.  Thus $\chi^m$ is constant.
Over characteristic zero a unit formal series with constant $m$th power
is constant, by comparing its first nonconstant homogeneous term; $m$
need only be nonzero.  Write $I$ for the given symmetry factor at $a$.
The normalization and symmetry identities are now the scalar equations
\[
 \chi=U_a(0)\kappa(a),\qquad
 \chi\kappa(-a)=U_a(-a),\qquad
 \kappa(-a)I=\kappa(a),\qquad
 U_a(0)I U_a(-a)=1.
\]
They give $\chi^2=1$.  Its order condition gives $\chi^m=1$, and $m$
odd implies $\chi=1$, as required.  The integral translation lifts,
their cube identities, and this odd-order normalization remain the given
theta data; no arbitrary action of a finite group is substituted.

Finally let $f_a=U_a t_a^*f$, let $\lambda\in K$, and let
$P\in K[X_{\tau i b}]$ be the original harmonic polynomial in $n$ columns,
homogeneous of total degree $d$.  For an integral series $F$ on $G^n$,
write
\[
 H_m(P;F)=\left[P(D)\left(e^{m\sum_bh_L(x_b)}F\right)\right]_{x=0}.
\]
This extends the previous column-product expression by linearity.
Suppose $\lambda\mu_*P$, obtained by the differential of $\mu$, has
integral coefficients and the required harmonicity.
The exact gauge identity is
\[
 e^{m\sum_bh_L(x_b)}\kappa\mu^*g
                   =\mu^*\left(e^{\sum_bh_L(x_b)}g\right).
\]
The chain rule for invariant derivatives and
\eqref{eq:formal-kernel-twisted-factor} consequently give
\[
       \lambda\sum_{a\in\ker\mu}H_m(P;f_a)=d_qj_p z,
                           \qquad \Delta^d z\in R,
\]
by Theorem~\ref{thm:discriminant-harmonic-clearing}.
For a general integral series $g$ on $G^n$, that theorem applies by
expanding its finite jet of order $d$ into integral products of monomials
in the individual columns.  Higher terms cannot affect a differential
value of order $d$.
In the CM tangent decomposition the unchanged Jacobian is
$j_p=\prod_\tau\tau(\det p)^{\operatorname{mult}(\tau)}$.

This proves the formal trace and divisibility step under the supplied
finite-height, splitting, regular-section, and symmetric-translation data.
Obtaining those data from the global model, passing from the formal kernel
to the full primary kernel, and comparing its cardinality with the
determinant ideals are separate remaining steps.
